\section{Experimental Setup}
\subsection{GLUE Benchmark}
The General Language Understanding Evaluation (GLUE) benchmark \cite{wang2018glue} consists of nine diverse NLP tasks (we use eight, excluding WNLI due to known issues). Each task evaluates different aspects of language understanding:
\begin{itemize}
    \item \textbf{CoLA}: Linguistic acceptability (binary classification, Matthews correlation).
    \item \textbf{SST-2}: Sentiment analysis (binary classification, accuracy).
    \item \textbf{MRPC}: Paraphrase detection (binary classification, accuracy/F1).
    \item \textbf{QQP}: Quora question pair duplication (binary, accuracy/F1).
    \item \textbf{QNLI}: Question answering entailment (binary, accuracy).
    \item \textbf{RTE}: Recognizing textual entailment (binary, accuracy).
    \item \textbf{MNLI}: Multi-genre natural language inference (3-class, accuracy).
    \item \textbf{STS-B}: Semantic textual similarity (regression, Pearson/Spearman).
\end{itemize}
We use the standard training/validation splits.
\subsection{Considerations for Small Datasets}
Several GLUE tasks (e.g., RTE, MRPC, STS‑B) have training sets of only a few thousand examples – STS‑B contains approximately 5.7k training samples and 1.5k validation samples. Fine‑tuning a 110M‑parameter model like BERT on such limited data risks overfitting. To mitigate this, we adopt:
\begin{itemize}
    \item A conservative learning rate ($2\times10^{-5}$) with linear warmup.
    \item Dropout of 0.1 (and 0.2 for smaller tasks).
    \item Early stopping based on validation loss.
    \item Evaluation after each epoch to monitor overfitting.
\end{itemize}
These practices are standard for low‑resource fine‑tuning and ensure that our SVD‑Lin comparisons are fair and robust.

\subsection{Model and Hyperparameters}
\begin{itemize}
    \item Base model: `bert-base-uncased` (110M parameters, 12 layers, 768 hidden size, 12 heads).
    \item Ranks tested: $r = 64, 128, 256$, plus full attention baseline ($r = N$).
    \item Batch size: 16 (training), 32 (evaluation).
    \item Learning rate: $2 \times 10^{-5}$ with linear warmup over 10\% of steps.
    \item Optimizer: AdamW ($\beta_1=0.9, \beta_2=0.999$, weight decay 0.01).
    \item Epochs: 3 for all tasks (except MNLI where we use 3 as well).
    \item Max sequence length: 128 tokens (truncation).
    \item Dropout: 0.1.
\end{itemize}

\subsection{Hardware and Software}
- **Machine**: Apple MacBook Pro M1 (2020), 16GB RAM, 8-core GPU.
- **Framework**: PyTorch 2.0.1 with MPS backend, Transformers 4.30.0, Datasets 2.14.0.
- **Reproducibility**: All seeds fixed to 42. Code available at \url{https://github.com/yourname/svd-lin}.

\subsection{Evaluation Metrics}
For each task and rank, we report the primary metric used by GLUE: Matthews correlation (CoLA), accuracy (most tasks), and Pearson correlation (STS-B). We also measure:
- **Training time per epoch** (seconds, averaged over 3 epochs).
- **Peak GPU memory** (using `torch.mps.driver\_allocated\_memory()` on MPS).
- **Inference latency** (ms per batch of 32 sequences).